\input ../common.tex

\graphicspath{{fig/}{../fig/}{../01/fig/}}
\usepackage[export]{adjustbox}

\institute[Nick Papior; DTU Compute]{\begin{tikzpicture}
      \node[shape=rectangle split,rectangle split parts=2,anchor=base] at (0,0)
      {SIESTA School};
    \end{tikzpicture}}

\date{21 November 2025}
\title{TranSiesta an intro to non-equilibrium Green function}
\author{Nick Papior}


\begin{document}

\begin{frame}
  \titlepage
\end{frame}

\begin{frame}
  \frametitle{Outline}
  \tableofcontents
\end{frame}

\showbondcurrentfalse
\showdensitymatrixweighingfalse
\showintromotivationfalse
\showintrotbtnterminalfalse
\shownegftransmissionfalse

\inputpath[../01]{introduction.tex}


\inputpath[../01]{green.tex}
\inputpath[../01]{se.tex}

\begin{frame}

  \vfill

  \centering
  \begin{minipage}{.7\linewidth}

  \begin{block}{Important}

    Understanding self-energies is like, really, really important.

  \end{block}

  \end{minipage}

  \vfill
\end{frame}

%\inputpath[../01]{negf.tex}
%\inputpath[../03]{contour.tex}

\inputpath[../03]{steps.tex}
%\inputpath[../03]{summary.tex}

\section{TranSiesta / TBtrans}
\begin{frame}
  \frametitle{TranSiesta / TBtrans}

  \begin{itemize}

    \item
      TranSiesta is a versatile and complex beast! There are \emph{many} pitfalls, mainly
      because of the \emph{self energy}.

      TranSiesta calculates the non-equilibrium \emph{density} and \emph{Hamiltonian}
      matrices.

    \item Start by converging parameters in plain Siesta calculations; then converge
      TranSiesta parameters, layers of electrodes (screening off the defect), energy
      points, etc.

    \item TBtrans (transmission function calculator) requires much finer $\mathbf k$ point
      sampling (recall slide \ref{integration} and \ref{green.bz})

      TBtrans uses a \emph{Hamiltonian} matrix (e.g. a non-equilibrium) and can calculate
      (projected) density of states, transmissions, etc.

    \item One can do $N_\mathfrak{e} \ge 1$ electrodes!
      \begin{center}
        \incg[height=8em,clip,trim=10cm 0 0 0]{ts-poisson-initial}
      \end{center}

    \item Completely removing periodic images possible:

      \begin{center}

        \incg[width=4cm]{rsse_self-energy}\hskip3em
        \incg[width=4cm]{rsse_replace}

      \end{center}


  \end{itemize}
  \doicite{Papior \etal: \doi{10.1103/PhysRevB.100.195417}}

\end{frame}

\section{Tutorials}
\begin{frame}
  \frametitle{Tutorials}

  \begin{center}

    There are additional tutorials covering NEGF using \texttt{sisl} on this web-page:

    \url{https://github.com/zerothi/ts-tbt-sisl-tutorial}

  \end{center}

\end{frame}

\end{document}
